\section{Numerical Simulations}\label{simulations}
We consider five second-order agents in a two-dimensional workspace.
The simulations illustrate the effect of tightening on transient
safety under the dynamic filter and the local exponential regime
established by Theorem~\ref{thm:closed_loop_LES}.
Their initial and target positions are
$[y_1(0),\cdots,y_5(0)] = \begin{bsmallmatrix} -3 & -3 & -3 & -3 & -3 \\ 0.002 & -1.605 & -0.55 & 0.55 & 1.605 \end{bsmallmatrix}$ and $[y_1^\star,\cdots,y_5^\star] = \begin{bsmallmatrix} 3 & 3 & 3 & 3 & 3 \\ 0 & -0.55 & -1.15 & 1.15 & 0.55 \end{bsmallmatrix}$.
All agents are initialized at rest, with
$
v_i(0)=\mathbf0_n
$
and
$
\mathbf y^i(0)=y(0)
$
for all $i\in\mathcal V$.
Complete simulation parameters and source code are available at
\url{https://github.com/SichenQian/safety-critical-gne}.

The cost of player $i$ is
$
J_i(y_i,y_{-i})
=
\frac{1}{2}\|y_i-p_i\|^2
+
\frac{0.03}{2}
\sum_{j\ne i}\|y_i-y_j\|^2
$ with
$[p_1,\cdots,p_5] = \begin{bsmallmatrix} 3 & 3 & 3 & 3 & 3 \\ 0 & -0.6325 & -1.3225 & 1.3225 & 0.6325 \end{bsmallmatrix}.$
The resulting pseudo-gradient is $1$-strongly monotone and
$1.15$-Lipschitz, and $y^\star$ is the unique unconstrained NE.
The sensing radius and inter-agent safety distance are
$d_c=3.2$ and $d_{ij}=0.345$, respectively.
The obstacles are centered at
$
o_1=(0.24,0)^\top$, $
o_{2,3}=(1.72,\pm1.1)^\top
$
with
$
r_{i,1}=0.625$ and $
r_{i,2}=r_{i,3}=0.315$ for $i\in\mathcal V.
$
We use the distance barriers
$
h_{ij}^{\mathrm{int}}
=
\|y_i-y_j\|^2-d_{ij}^2$ and
$
h_{i,k}^{\mathrm{obs}}
=
\|y_i-o_k\|^2-r_{i,k}^2$, together with the matrix-tree certificate in
Remark~\ref{tree}, with $\underline{\lambda}_\chi=0.2$. For numerical conditioning, the matrix-tree certificate is normalized
by $512$. The physical collision, obstacle, and connectivity ECBF
terms and their tightening terms are multiplied by $0.06$, $0.06$,
and $1$, respectively; the mismatch coefficients retain their
original definitions.
The parameters in \eqref{stacked_closed_loop_flow} are set as
$K_1 = 14.5$, $K_2 = 8.8$, $c = 40$, $\gamma = 45$, $\xi = 1.4$, $\epsilon = 2\times10^{-7}$, $\Delta_{ij}^{\mathrm{int}} = 1.2$, $\Delta_{i,k}^{\mathrm{obs}} = 1.8$, $\Delta^{\mathrm{c}} = 2.5$, $\nu = 0.1$.
Moreover, $\theta_i=50$ and $\delta_{p,i}=10^{-3}$ for all
$i\in\mathcal V$. For this game, the extended pseudo-gradient has
$\tilde l\simeq1.1216$, and hence $c=40$ and $\gamma=45$
satisfy \eqref{cond_c}--\eqref{cond_gamma}. For each variant, the fast variables are initialized at a
primal--dual solution of its corresponding initial regularized QP.
Solver settings are provided in the code repository. The simulation focuses on the safety-filtering effect and the local
exponential regime, while the sufficient entrance condition in
Corollary~\ref{cor:regional_to_local} is not invoked for this experiment.

To isolate the effect of the tightening margins, we compare the
proposed flow with an untightened ablation using identical initial
physical states, local decision estimates, and algorithm parameters,
but with
$
\Delta_{ij}^{\mathrm{int}}
=
\Delta_{i,k}^{\mathrm{obs}}
=
\Delta^{\mathrm c}
=
0.
$
Figs.~\ref{fig:traj-tightened} and~\ref{fig:traj-untightened}
show the corresponding trajectories. Under the tightened flow, all
agents avoid the obstacles and approach the prescribed target
configuration, whereas the highlighted interaction between agent~$1$
and obstacle~$1$ has a small negative sampled clearance in the untightened case.
Fig.~\ref{fig:safety-margins} compares the transient safety
quantities. All three remain positive in the tightened run. The
untightened collision and obstacle minima are $-4.24\times10^{-7}$
and $-9.73\times10^{-7}$, respectively, whereas its matrix-tree
minimum is positive, $8.33\times10^{-6}$. These signs persist under
tighter integration tolerances.

The distance-dependent edge $(2,5)$ activates naturally at
$t\simeq0.049766$~s in the tightened run, increasing the number of
edges from nine to ten. Over the saved states, the minimum
$\lambda_2(L_\chi(y))$ is $1.93982$ in both variants, and the
tightened matrix-tree minimum is $0.028725$.
At 30,001 saved times, with denser sampling during the initial and
active-filter transients, the maximum regularized-QP tracking error
is $0.0207623$. The connectivity discrepancy maximized over
$\|u\|\le35$ at each sampled state is at most $0.903881$.
Using $\bar a^{\mathrm{int}}=6.44$,
$\bar a^{\mathrm{obs}}=11.14852$, and
$\bar a^{\mathrm c}=1.85130$, the three tightening margins are
$1.01091$, $1.56853$, and $1.51017$.
All sampled QPs satisfy feasibility and KKT checks within
$10^{-8}$, and the original physical ECBF residuals are nonpositive
along the tightened run. These trajectory-based numerical
diagnostics support the tightening mechanism; they do not establish
the uniform regional bounds required by Theorem~\ref{thm:regional_safety}.
Finally, Fig.~\ref{fig:state-convergence} shows convergence of the
position components to $y^\star$.
Fig.~\ref{fig:local-exp} uses the same tightened trajectory and plots
the aggregate error
$
\sqrt{\|e_\vartheta(t)\|^2+\|e_f(t)\|^2}
$
over $t\in[850,1050]$.
Its approximately linear semilogarithmic decay is consistent with the
local exponential convergence established in
Theorem~\ref{thm:closed_loop_LES}.
\begin{figure}[htbp]
	\centering
	\subfloat[Tightened flow; inset: positive obstacle clearance.]{
		\includegraphics[width=\linewidth]{figures/my_trajectory.pdf}\label{fig:traj-tightened}}
	\hfill
\subfloat[Untightened flow; inset: obstacle penetration.]{
		\includegraphics[width=\linewidth]{figures/my_trajectory_2.pdf}\label{fig:traj-untightened}}
	\caption{Planar trajectories under the tightened and untightened flows.}
	\label{fig:trajectory-comparison}
\end{figure}
\begin{figure}[htbp]
	\centering
	\subfloat[Inter-agent collision margin]{
		\includegraphics[width=\linewidth]{figures/inter_agent_avoidance.pdf}
		\label{fig:inter_agent_avoidance}
	}
\\
	\subfloat[Obstacle-avoidance margin]{
		\includegraphics[width=\linewidth]{figures/static_obstacle_avoidance.pdf}
		\label{fig:static_obstacle_avoidance}
	}\\
	\subfloat[Matrix-tree connectivity certificate]{
		\includegraphics[width=\linewidth]{figures/connectivity_preservation.pdf}
		\label{fig:connectivity_preservation}
	}
\caption{Safety quantities under the tightened and untightened flows.} \label{fig:safety-margins}
\end{figure}
\begin{figure}[htbp]
	\centering
\subfloat[Position-component convergence.]{
		\includegraphics[width=\linewidth]{figures/HB_flow_hebing.pdf}
		\label{fig:state-convergence}
	}
	\\
	\subfloat[Local decay of the closed-loop error along the same tightened
		trajectory.]{
		\includegraphics[width=\linewidth]{figures/exp_2.pdf}
		\label{fig:local-exp}
	}
\caption{Convergence behavior under the tightened flow.}
	\label{fig:convergence}
\end{figure}
